% Originally: content/week-12.tex, \section{Pullback of forms} (Corsin Week 12)


\section{Pullback of forms}
\label{sec:pullback_of_forms}

% Source: Corsin Nick/Class Notes/Week 12.pdf, p. 5
\begin{definition}[Pullback of a form]
\label{def:pullback_of_forms}
Let $U \subseteq \mathbb{R}^m$, $V \subseteq \mathbb{R}^n$ be open, $F : U \to V$ a smooth
map, and $\omega \in \Omega^p(V)$ a differential $p$-form on $V$. The \newterm{pullback}
(\germanterm{Zur\"uckziehung}) of $\omega$ under $F$ is the differential $p$-form on $U$
defined by
\[(F^*\omega)_x(v_1,\ldots,v_p) := \omega_{F(x)}\big(DF_x(v_1),\ldots,DF_x(v_p)\big).\]
\end{definition}

\begin{ainote}
Corsin also draws a commutative diagram $TU \xrightarrow{dF} TV \xrightarrow{\omega} \mathbb{R}$
over $U \xrightarrow{F} V$, with tangent-bundle projections $\pi_U, \pi_V$, and labels it
\qt{not exam relevant}. Omitted here for that reason; the formula above is the operative
content.
\end{ainote}

The pullback satisfies two properties that make it the right tool for transporting forms
along smooth maps:
\begin{enumerate}[label=\textbf{(\arabic*)}]
  \item $F^*(\omega\wedge\theta) = F^*\omega \wedge F^*\theta$;
  \item $F^*(d\omega) = d(F^*\omega)$ (the pullback commutes with the exterior derivative).
\end{enumerate}

\begin{ainote}
This is precisely the operation used, informally, in the reparametrization-invariance
motivation for \cref{def:antisymmetric_k_linear_form}: writing $\phi^*\omega$ and
$(\phi\circ\Psi)^*\omega$ for the two sides of that calculation makes it read as $\int_D
\phi^*\omega = \int_V (\phi\circ\Psi)^*\omega = \int_V \Psi^*(\phi^*\omega)$, i.e.\ property (1)
applied to $\phi^*\omega$ and the fact that pullback along a composition is the composition of
pullbacks.
\end{ainote}
